\section{Conclusion}
We introduced \dysta{}, a method that improves efficiency and multi-frame integration in VLA models via static–dynamic disentanglement. By reusing temporally persistent static tokens and selectively refreshing them through a learnable recache gate, \dysta{} enables compact multi-frame contexts and efficient inference. Experiments show that our approach substantially improves performance on temporally dependent tasks while providing meaningful speedups over strong baselines. These results suggest that explicitly modeling temporal persistence is a promising direction for scalable and practical VLA systems. \noindent\textbf{Limitations}. We realize our approach based on pretrained VLAs. Such a strategy may not fully unleash the performance of the model. Future work could focus on pretraining VLAs with our architecture from scratch.